% Einheit 3 von 4 – Gerade und Ebene (bank/lagebeziehungen/e3.jsonl, zusammenbau v0.1)
%% TODO Zeitmarke Einheit 3: Marken-Zeile nicht lesbar
%% TODO Zweigzeile Teil 1: was hier gelernt wird (Satz in Schülersprache aus dem Einheitstitel) – Einheit 3
\einheitenkopf[e3]{Einheit 3 von 4 · Gerade und Ebene}
\zweigzeile{Abitur GK}
\setcounter{aufgabe}{13}

%% TODO Titel: Ich-kann-Satz fehlt in der Bank (Platzhalter: Kettenname) – Nr. 14
%% TODO Anweisung: ein Satz über der ersten Teilaufgabe fehlt in der Bank – Nr. 14
\begin{aufgabe}{Gerade und Ebene}
%% TODO \rechenplatz{n}: Schrittzahl des Grundfalls fehlt in der Bank (nur bei fester Schreibform, 2.3 b)
\begin{teile}
\teil Der allgemeine Punkt von $g$ in die Gleichung von $E$ eingesetzt ergibt $4 + 3t - 3t = 4$. Wie liegt $g$ zu $E$? \\ \kreuz{liegt in der Ebene} \\ \kreuz{echt parallel} \\ \kreuz{genau ein Schnittpunkt}
\teil Weise nach, dass $g: \vec{x} = (1 | 2 | 0) + t \cdot (1 | -1 | 1)$ in $E: x + 2y + z = 5$ liegt.
\teil $g: \vec{x} = (1 | 0 | 2) + t \cdot (2 | 1 | 1)$ und $E: x - 2y + z = 1$. Entscheide mit der Merkregel, wie $g$ zu $E$ liegt.
\teil $P(1 | 3 | 8)$ und $Q_t(5 | t | 6)$ mit reellem $t$. Gibt es ein $t$, für das die Gerade $PQ_t$ parallel zur xy-Ebene verläuft? Begründe. \hfill (Abitur 2024 GK)
\teil Der Punkt $S(4 | 3 | 5)$ wird um die Achse $a: \vec{x} = (0 | 0 | 2) + t \cdot (1 | 2 | 0)$ gedreht. $L: x + 2y = 10$; $H$ ist der Schnittpunkt von $a$ und $L$. Begründe, dass $S$ auf einem Kreis in $L$ mit dem Mittelpunkt $H$ läuft, und bestimme $H$. \hfill (Abitur 2026 LK)
\teil Begründe, dass es unendlich viele Ebenen gibt, die keinen Punkt der Form $(a | a | 0)$ enthalten. \hfill (Abitur 2019 LK)
\teil $E: 2x - z = 1$ und für jede Zahl $a$ die Gerade $g_a: \vec{x} = (1 | 0 | a) + r \cdot (a | 2 | 2)$. Untersuche die Lage von $g_a$ zu $E$ in Abhängigkeit von $a$ und gib gegebenenfalls den Schnittpunkt an. \hfill (Abitur 2024 LK)
\end{teile}
\end{aufgabe}

%% TODO Titel: Ich-kann-Satz fehlt in der Bank (Platzhalter: Kettenname) – Nr. 15
%% TODO Anweisung: ein Satz über der ersten Teilaufgabe fehlt in der Bank – Nr. 15
\begin{aufgabe}{Gerade und Ebene – Fehler finden, Begründen}
\begin{teile}
\teil Ich finde den Fehler: $g: \vec{x} = (1 | 2 | 0) + t \cdot (1 | 0 | 1)$, $E: 2x - y + z = 0$. Nils setzt den Stützpunkt ein, $2 - 2 + 0 = 0$, und schreibt: $g$ liegt in $E$.
\teil Begründe: Fällt beim Einsetzen des allgemeinen Geradenpunkts der Parameter heraus und die Gleichung stimmt, liegt die ganze Gerade in der Ebene.
\end{teile}
\end{aufgabe}
